\chapter{Menopause}
\index{Menopause}

\section*{Definition and Overview}
Menopause is a natural biological milestone in a woman's life, representing the permanent cessation of menstruation due to the loss of ovarian follicular activity and subsequent oestrogen depletion. Clinically, natural menopause is diagnosed retrospectively after 12 consecutive months of amenorrhoea in the absence of other physiological or pathological causes. The transition period preceding menopause is termed perimenopause (or the menopausal transition), during which fluctuating hormone levels lead to irregular menstrual cycles and various physical and psychological symptoms. The postmenopausal phase begins immediately after the 12-month milestone has been reached. Menopause can also occur prematurely or be induced medically (via chemotherapy or pelvic radiation) or surgically (via bilateral oophorectomy).

\section*{Epidemiology}
The median age of natural menopause globally and in many countries is approximately 51 years, with the majority of women experiencing menopause between the ages of 45 and 55. Menopause occurring between the ages of 40 and 45 is termed early menopause, whereas menopause occurring before the age of 40 is classified as premature ovarian insufficiency (POI), which affects approximately 1\% of the female population. The perimenopausal transition typically begins in a woman's mid-to-late 40s and lasts between 4 and 8 years. Virtually all women who reach mid-life will navigate this transition, although the duration, severity, and impact of menopausal symptoms show significant individual, geographic, and cultural variation.

\section*{Aetiology and Pathophysiology}
The fundamental pathophysiology of menopause is the age-related depletion of the ovarian follicular pool:
\begin{itemize}
    \item \textbf{Follicular exhaustion:} As a woman ages, the number of primordial follicles in the ovaries declines. Eventually, the remaining follicles fail to respond to gonadotropins, leading to a profound decrease in the production of oestradiol and progesterone.
    \item \textbf{Gonadotropin rise:} In response to the loss of negative feedback from oestradiol and inhibin, the anterior pituitary gland secretes markedly elevated levels of follicle-stimulating hormone (FSH) and luteinising hormone (LH).
    \item \textbf{Perimenopausal fluctuations:} During perimenopause, oestrogen levels do not decline linearly but fluctuate unpredictably, driving menstrual cycle irregularity and acute vasomotor episodes.
    \item \textbf{Surgical oophorectomy:} Bilateral removal of the ovaries causes an immediate and severe drop in systemic oestrogen and androgen levels, resulting in abrupt and often intense menopausal symptoms.
    \item \textbf{Gonadotoxicity:} Cancer treatments, such as alkylating chemotherapy or pelvic radiotherapy, damage the sensitive ovarian follicular structure, causing temporary or permanent ovarian failure.
    \item \textbf{Premature Ovarian Insufficiency (POI):} Ovarian failure before age 40, which may be idiopathic, autoimmune (e.g. associated with adrenalitis), or genetic (e.g. Turner syndrome mosaicism, Fragile X premutation).
\end{itemize}

\section*{Clinical Features}
The clinical presentation of menopause spans menstrual changes, vasomotor instability, genitourinary atrophy, and cognitive fluctuations:
\begin{itemize}
    \item \textbf{Menstrual cycle alterations:} Irregular menstrual bleeding patterns, changes in flow volume (exceptionally heavy or light), and fluctuating cycle lengths during perimenopause, leading to amenorrhoea.
    \item \textbf{Vasomotor symptoms:} Hot flushes (sudden, intense sensations of heat spreading over the face, neck, and chest) and night sweats, which are frequently accompanied by transient palpitations and anxiety.
    \item \textbf{Genitourinary Syndrome of Menopause (GSM):} Vulvovaginal dryness, pruritus, dyspareunia (painful intercourse), urinary urgency, dysuria, and recurrent urinary tract infections (UTIs) caused by oestrogen-deficient thinning of the urogenital epithelium.
    \item \textbf{Sleep disturbances:} Insomnia, frequent nocturnal awakenings, and poor sleep quality, which are often triggered by night sweats.
    \item \textbf{Neuropsychiatric changes:} Mood swings, irritability, anxiety, mild depressive symptoms, and cognitive changes often described as ``brain fog'' (forgetfulness and poor concentration).
    \item \textbf{Somatic changes:} Diffuse joint and muscle stiffness (arthralgia), breast tenderness, dry skin, hair thinning, and a redistribution of body fat leading to increased abdominal adiposity.
\end{itemize}

\section*{Differential Diagnosis}
Because the symptoms of perimenopause and menopause overlap with other medical conditions, clinicians must rule out:
\begin{itemize}
    \item \textbf{Thyroid disease:} Hypothyroidism (presenting with fatigue, weight gain, and irregular menses) and hyperthyroidism (causing heat intolerance, sweating, and palpitations), distinguished by a thyroid function panel.
    \item \textbf{Pregnancy:} In any perimenopausal woman presenting with amenorrhoea or irregular bleeding, pregnancy must be actively excluded with a beta-hCG test.
    \item \textbf{Hyperprolactinaemia:} Elevated prolactin levels due to a pituitary adenoma or medications, presenting with amenorrhoea and galactorrhoea.
    \item \textbf{Primary psychiatric disorders:} Major depressive disorder or generalized anxiety disorder, which present with chronic mood changes but lack the characteristic vasomotor and genitourinary symptoms.
    \item \textbf{Systemic malignancies:} Conditions such as lymphoma or carcinoid syndrome, which can cause severe night sweats and flushing, distinguished by constitutional symptoms (unintentional weight loss, fever) and diagnostic investigations.
\end{itemize}

\section*{When to Seek Medical Care}
Women should consult a general practitioner to discuss symptom management if menopausal symptoms are impairing their quality of life, or if they experience symptoms of menopause before the age of 45.

\textbf{Contact your local emergency services for:}
\begin{itemize}
    \item Sudden, severe, or crushing chest pain that may radiate to the jaw or arm.
    \item Sudden weakness or numbness in the face, arm, or leg, or sudden difficulty speaking (signs of stroke).
    \item Sudden shortness of breath or loss of consciousness.
\end{itemize}
Any vaginal bleeding occurring after 12 months of amenorrhoea (postmenopausal bleeding) is abnormal and requires urgent investigation by a doctor.

\section*{Diagnosis and Investigations}
In women over the age of 45 with typical symptoms, menopause is diagnosed clinically. In other cases, investigations are required:
\begin{itemize}
    \item \textbf{Clinical evaluation:} Detailed history of menstrual patterns, symptoms, and medical history.
    \item \textbf{Follicle-stimulating hormone (FSH) levels:} Measurement of serum FSH is not routinely required in women over 45, but is indicated in women under 40 suspected of POI (diagnosed when FSH is $>30$ IU/L on two occasions 4--6 weeks apart).
    \item \textbf{Thyroid-stimulating hormone (TSH):} Assayed to rule out thyroid dysfunction in patients presenting with hot flushes and menstrual changes.
    \item \textbf{Transvaginal ultrasound (TVUS):} Indicated to assess endometrial thickness and rule out endometrial hyperplasia or malignancy in women presenting with postmenopausal bleeding or abnormal perimenopausal bleeding.
    \item \textbf{Dual-energy X-ray absorptiometry (DEXA) scan:} Recommended to establish baseline bone mineral density (BMD) in women with early menopause, POI, or other risk factors for osteoporosis.
\end{itemize}

\section*{Management}
Management focuses on symptom relief, health education, and the prevention of long-term complications:
\begin{itemize}
    \item \textbf{Menopausal Hormone Therapy (MHT):} The most effective treatment for vasomotor symptoms. It involves systemic oestrogen-only therapy for women who have undergone a hysterectomy, or combined oestrogen-progestogen therapy for women with an intact uterus to prevent endometrial hyperplasia.
    \item \textbf{Local vaginal therapies:} Low-dose topical oestrogen (creams, pessaries, or rings) to manage GSM, providing localised relief with minimal systemic absorption.
    \item \textbf{Non-hormonal pharmacotherapy:} SSRIs or SNRIs (e.g. venlafaxine, desvenlafaxine), gabapentin, or clonidine to alleviate vasomotor symptoms in women with contraindications to MHT (such as a history of oestrogen-receptor-positive breast cancer).
    \item \textbf{Cognitive Behavioural Therapy (CBT):} Evidence-based psychological therapy to reduce the perceived severity of hot flushes, improve sleep hygiene, and manage mood fluctuations.
    \item \textbf{Lifestyle interventions:} Wearing layered clothing, maintaining a cool environment, engaging in regular weight-bearing and resistance exercise, and reducing triggers (alcohol, caffeine, spicy foods).
    \item \textbf{Bone-specific therapies:} Adequate intake of calcium and vitamin D, and bisphosphonates or denosumab in women with diagnosed osteoporosis.
\end{itemize}

\section*{Complications}
The long-term physiological consequences of postmenopausal oestrogen deficiency include:
\begin{itemize}
    \item \textbf{Osteoporosis:} Rapid acceleration of bone resorption leads to decreased bone mineral density and an elevated risk of fragility fractures (particularly of the hip, vertebrae, and wrist).
    \item \textbf{Cardiovascular disease:} Loss of the cardioprotective effects of oestrogen on lipid profiles and vascular endothelial function, leading to an increased risk of coronary artery disease and stroke.
    \item \textbf{Genitourinary dysfunction:} Progressive vaginal stenosis, severe dyspareunia, pelvic organ prolapse, and stress or urge urinary incontinence.
    \item \textbf{Metabolic shifts:} Increased visceral fat deposition, adverse lipid changes (elevated LDL and triglycerides, decreased HDL), and an increased risk of developing type 2 diabetes.
    \item \textbf{Cognitive changes:} Increased subjective memory complaints, although the exact relationship between menopause and long-term dementia risk remains a subject of ongoing research.
\end{itemize}

\section*{Prognosis and Outcomes}
Menopause is a natural transition, and the long-term prognosis is excellent. Vasomotor symptoms typically peak during the late transition and early postmenopausal years, resolving spontaneously within 4 to 7 years in most women, though they can persist for a decade or longer in some. Genitourinary symptoms (GSM), however, do not resolve spontaneously and tend to progress without intervention. With the judicious use of Menopausal Hormone Therapy, lifestyle modifications, and modern supportive treatments, most women successfully manage their symptoms and lead active, healthy lives post-menopause.

\section*{Prevention and Self-Care}
Self-care and preventative health measures are vital to protect bone and cardiovascular health:
\begin{itemize}
    \item \textbf{Smoking cessation:} Smoking increases the clearance of oestrogen, worsens vasomotor symptoms, and accelerates bone loss, making quitting a priority.
    \item \textbf{Regular weight-bearing exercise:} Engage in regular walking, jogging, weight training, or resistance exercises to stimulate bone remodeling and maintain muscle mass.
    \item \textbf{Nutrition for bone health:} Ensure a diet rich in calcium (dairy products, fortified plant milks, leafy green vegetables) and maintain adequate vitamin D through sensible sun exposure or supplements.
    \item \textbf{Identify vasomotor triggers:} Keep a diary to pinpoint triggers like caffeine, alcohol, spicy foods, or heated blankets, and adjust habits accordingly.
    \item \textbf{Vaginal moisturisers:} Regular use of non-hormonal, water-based vaginal moisturisers and lubricants to ease dryness and maintain tissue elasticity.
\end{itemize}

\section*{Sources and Further Reading}
The following reputable organisations provide further information on menopause:
\begin{itemize}
    \item Australasian Menopause Society (AMS). \textit{Clinical resources, patient information, and guides on hormone therapy.} https://www.menopause.org.au/
    \item Jean Hailes for Women's Health. \textit{Comprehensive health information on midlife health, perimenopause, and menopause.} https://www.jeanhailes.org.au/
    \item NHS. \textit{Overview of menopause symptoms, diagnosis, and management options.} https://www.nhs.uk/
    \item Mayo Clinic. \textit{Clinical advice on menopause transition, symptoms, and therapies.} https://www.mayoclinic.org/
    \item The Menopause Society (formerly NAMS). \textit{Education and guidelines on menopausal health and treatment.} https://www.menopause.org/
\end{itemize}


\section*{Definition and Overview}
Menopause is the permanent cessation of menstruation, confirmed retrospectively after 12 consecutive months without a period. It represents the natural end of a woman's reproductive years, when the ovaries stop releasing eggs and markedly reduce their production of oestrogen and progesterone. The years leading up to the final period, during which hormone levels fluctuate and symptoms often begin, are known as perimenopause or the menopausal transition. Menopause can also be induced by surgery, such as removal of both ovaries, or by medical treatments that damage ovarian function. This chapter provides an overview of the physiological changes, common symptoms, and the treatment and self-care options available to manage this stage of life.

\section*{Epidemiology}
Menopause is a universal biological event that affects every woman who lives to midlife. The average age of natural menopause is around 51 years in many countries and comparable populations, with most women experiencing the transition between their mid-40s and mid-50s. Approximately 1 in 10 women experience premature menopause before the age of 40, or early menopause between 40 and 45 years, which may be natural or caused by surgery, chemotherapy, or other factors. Up to three-quarters of women report vasomotor symptoms such as hot flushes and night sweats during the transition, and for roughly one in four these symptoms are severe enough to interfere with daily functioning. With increasing female life expectancy, many women now spend a substantial proportion of their lives in the post-menopausal state, making the long-term health effects of oestrogen deficiency an important public health concern.

\section*{Aetiology and Pathophysiology}
The central event in menopause is the age-related depletion of ovarian follicles, the structures that produce eggs and the steroid hormones oestrogen and progesterone. As follicular numbers decline, ovulation becomes erratic and eventually ceases, and circulating oestrogen levels fall to a low, stable plateau.

\begin{itemize}
    \item \textbf{Natural ageing:} the ovaries' finite supply of eggs is gradually exhausted, reducing hormone output and leading to irregular cycles that ultimately stop
    \item \textbf{Surgical menopause:} bilateral oophorectomy causes an abrupt and often more severe onset of symptoms because hormone levels fall suddenly
    \item \textbf{Medical treatments:} chemotherapy, pelvic radiotherapy, and some endocrine therapies can damage or suppress ovarian function, sometimes permanently
    \item \textbf{Premature ovarian insufficiency:} ovarian failure before age 40, which may be autoimmune, genetic (for example in Turner syndrome), or of unknown cause
    \item \textbf{Effects of oestrogen deficiency:} low oestrogen affects the hypothalamus, destabilising temperature regulation and causing vasomotor flushes; it also thins the vaginal epithelium, reduces bone density, and influences mood and sleep pathways
\end{itemize}

\section*{Clinical Features}
The symptoms of menopause vary widely between women in type and severity. The hallmark symptoms are vasomotor: hot flushes and night sweats, which can disrupt sleep and lead to significant daytime fatigue.

\begin{itemize}
    \item Hot flushes and night sweats, the most common and characteristic symptoms
    \item Sleep disturbance, insomnia, and resulting tiredness and irritability
    \item Vaginal dryness, itching, and discomfort, often with dyspareunia and reduced libido
    \item Mood changes, anxiety, irritability, and low mood
    \item Joint and muscle aches, headaches, and a sensation of "brain fog" or poor concentration
    \item Menstrual irregularity, with periods becoming longer, shorter, lighter, or heavier before stopping
    \item Urinary symptoms such as urgency, frequency, and recurrent urinary tract infections
    \item Skin and hair changes, including dryness and thinning
\end{itemize}

\section*{Differential Diagnosis}
Although menopause is a normal physiological event, its symptoms overlap with several medical conditions that should be considered, particularly when they are atypical or severe.

\begin{itemize}
    \item \textbf{Thyroid disease:} hypothyroidism and hyperthyroidism can both cause fatigue, mood disturbance, weight change, and menstrual irregularity
    \item \textbf{Anaemia:} iron deficiency may produce tiredness, breathlessness, and palpitations
    \item \textbf{Depression or anxiety:} primary mood disorders can occur independently of the menopause
    \item \textbf{Chronic fatigue syndrome} and other causes of persistent exhaustion
    \item \textbf{Pregnancy:} possible in women who are still ovulating intermittently
    \item \textbf{Structural gynaecological causes of bleeding:} uterine fibroids, polyps, or rarely endometrial cancer, especially with heavy or post-menopausal bleeding
\end{itemize}

\section*{When to Seek Medical Care}
Women should consult a doctor if menopausal symptoms are troublesome or affecting quality of life, if periods become very heavy or unpredictable, if bleeding occurs more than 12 months after the last period, or if they wish to discuss the risks and benefits of hormone therapy.

\textbf{Contact your local emergency services for:}
\begin{itemize}
    \item Very heavy vaginal bleeding that soaks pads or tampons within an hour, especially with faintness or dizziness
    \item Heavy or painful bleeding after a long gap since the last period
    \item Sudden, severe abdominal or pelvic pain
    \item Chest pain, sudden breathlessness, or painful swelling of one leg, which may indicate a blood clot
    \item Signs of stroke: sudden facial drooping, arm weakness, or difficulty speaking
\end{itemize}

\section*{Diagnosis and Investigations}
In most cases the diagnosis of menopause is made clinically from the history of amenorrhoea and typical symptoms, and no tests are required for women over the age of 45. Investigations are used mainly to exclude other causes or to investigate abnormal bleeding.

\begin{itemize}
    \item \textbf{History:} menstrual pattern, symptom diary, timing of symptoms, and impact on sleep and daily life
    \item \textbf{Examination:} blood pressure, body mass index, and pelvic examination when bleeding symptoms or gynaecological concerns are present
    \item \textbf{Blood tests:} follicle-stimulating hormone (FSH) measurement may help confirm menopause in younger women or those in whom the diagnosis is unclear; it is not generally needed over age 45
    \item \textbf{Thyroid function tests and full blood count:} to exclude thyroid disease and anaemia as causes of symptoms
    \item \textbf{Pelvic ultrasound:} for heavy or irregular bleeding, to assess the endometrium, uterus, and ovaries
    \item \textbf{Endometrial sampling:} in post-menopausal bleeding, to exclude endometrial hyperplasia or cancer
\end{itemize}

\section*{Management}
Management is individualised and combines lifestyle measures with pharmacological treatment where symptoms warrant it. Treatment should be reviewed regularly as needs change.

\begin{itemize}
    \item \textbf{Lifestyle measures:} layered clothing, cooling techniques, regular exercise, and avoidance of alcohol, caffeine, and spicy food before sleep
    \item \textbf{Hormone replacement therapy (HRT):} the most effective treatment for vasomotor symptoms; oestrogen is used alone in women without a uterus and combined with progestogen in those with one, to protect the endometrium
    \item \textbf{Vaginal oestrogen:} low-dose topical therapy for genitourinary symptoms, with minimal systemic absorption
    \item \textbf{Non-hormonal medication:} some antidepressants and other agents can reduce hot flushes in women who cannot take or prefer to avoid HRT
    \item \textbf{Talking therapies:} cognitive behavioural therapy and mindfulness have evidence for improving vasomotor symptoms, sleep, and mood
    \item \textbf{Bone health:} adequate calcium and vitamin D, weight-bearing exercise, and bone density assessment where risk factors exist
    \item \textbf{Complementary approaches:} some women find benefit from herbal preparations, though evidence of effectiveness is limited and products are not tightly regulated
\end{itemize}

\section*{Complications}
\begin{itemize}
    \item Osteoporosis and increased fracture risk resulting from accelerated bone loss in the post-menopausal years
    \item Increased cardiovascular risk, including coronary heart disease and stroke, after menopause
    \item Genitourinary symptoms of menopause that affect intimacy, continence, and quality of life
    \item Chronic sleep disturbance and adverse effects on mood, work performance, and relationships
    \item Potential HRT side effects, including a small increased risk of venous thromboembolism, stroke, breast cancer, and gallbladder disease in some women, which must be weighed against benefits in shared decision-making
\end{itemize}

\section*{Prognosis and Outcomes}
Most women find that the most troublesome menopausal symptoms settle over time. Vasomotor symptoms typically improve within 2--5 years of the final period, although a minority of women continue to experience flushes for a decade or longer. HRT is highly effective at relieving vasomotor and many other symptoms while it is used. The long-term consequences of oestrogen deficiency, particularly bone loss and cardiovascular risk, persist after menopause and respond to lifestyle measures and, where appropriate, pharmacological bone and vascular protection. Women who experience premature or early menopause face greater cumulative health risks and are generally advised to consider hormone therapy until at least the natural age of menopause, with the balance of risks and benefits discussed individually.

\section*{Prevention and Self-Care}
\begin{itemize}
    \item Discuss symptoms openly with a doctor rather than accepting them as an unavoidable burden
    \item Maintain regular physical activity, including weight-bearing and resistance exercise
    \item Eat a balanced diet rich in calcium and vitamin D, and consider supplements if intake is inadequate
    \item Avoid smoking and limit alcohol, both of which can trigger flushes and worsen bone and cardiovascular health
    \item Use practical cooling measures: fans, cool drinks, layered clothing, and a cool bedroom
    \item Protect sleep with a consistent bedtime routine and treat any sleep disorder early
    \item Seek support for low mood, anxiety, or relationship strain, including talking therapies
    \item Consider bone density screening if there are risk factors for osteoporosis
\end{itemize}

\section*{Sources and Further Reading}
The following reputable sources provide further information for healthcare students and patients seeking reliable, current advice about the menopause and its management.

\begin{itemize}
    \item NHS. \textit{Menopause: symptoms, diagnosis, and treatment options.} https://www.nhs.uk/conditions/menopause/
    \item Mayo Clinic. \textit{Patient guide to menopause and hormone therapy.} https://www.mayoclinic.org/diseases-conditions/menopause/
    \item MSD Manual. \textit{Menopause: professional overview and management.} https://www.msdmanuals.com/
    \item MedlinePlus. \textit{Menopause health information.} https://medlineplus.gov/menopause.html
    \item World Health Organization. \textit{Menopause and women's health.} https://www.who.int/
    \item NICE. \textit{Menopause: diagnosis and management guidance.} https://www.nice.org.uk/guidance/ng23
\end{itemize}
